\section{Can Scalable Oversight Correct Rater Bias?}
\label{sec:11.4}

\textbf{Unresolved issue.}
Preference optimization can teach a model to reproduce a rater's systematic errors, including the preference for agreeable answers over accurate ones \cite{sharma2023towards}. It remains unresolved whether weak-to-strong oversight allows a more capable model to exceed a supervisor's measured bias or merely reproduces that bias at greater scale.

\textbf{Test.}
Measure prospective raters' susceptibility to a documented bias, then use raters with different measured bias levels to supervise otherwise matched stronger models. Evaluate the resulting models on questions for which the raters' beliefs and independently established answers diverge. Compare whether each model tracks the ground truth, the rater's competence, or the rater's errors.

\textbf{Evidence against the proposal.}
If stronger models reproduce supervisors' errors in proportion to the supervisors' measured bias, latent model competence does not correct a defective training signal. Scalable oversight would then be unable to support the proposal wherever the relevant human judgments are systematically distorted.

\textbf{Required access or authority.}
The experiment requires rater-level bias measurements, training records linking ratings to resulting models, and an independently established answer key. Researchers outside the training organization must control the evaluation and retain the right to publish results that expose failures in the preference pipeline.
